% ==========================================
%  content/manual-ch01-intro.tex  —  第一章 绪论
% ==========================================
\documentclass[../main-manual.tex]{subfiles}
\begin{document}

\cleardoublepage
\thispagestyle{empty}
\zihao{-4} \songti
\setlength{\baselineskip}{24pt}
\pagenumbering{arabic}
\RaggedRight
\setlength{\parindent}{2em}
\raggedbottom

\chapter{绪论} \label{ch:intro}

\section{\TeX{} 与 \LaTeX{} 的历史与意义}

\subsection{\TeX{} 的诞生}

1978年，斯坦福大学计算机科学教授 Donald E. Knuth 在修订其经典著作《The Art of Computer Programming》第二卷时，对出版商提供的排版质量深感不满。当时的照相排版技术无法令人满意地渲染数学公式，于是 Knuth 做了一个影响至今的决定——暂停写作，先开发一套能够高质量排版数学内容的计算机系统。

这一决定催生了 \TeX{}（读作 "tech"，取自希腊语 $\tau\varepsilon\chi\nu\eta$，意为"艺术"与"技术"）。Knuth 不仅实现了断行算法、数学公式渲染、连字符处理等核心功能，还将整个系统的源代码以 literate programming 的方式公开\cite{Knuth1984}，开创了自由软件的先河。\TeX{} 的版本号收敛于 $\pi$（当前为 3.141592653），象征着系统的稳定性与永恒性。

\subsection{\LaTeX{} 的出现与发展}

尽管 \TeX{} 功能强大，但其底层命令较为原始，普通用户需要编写大量样板代码来完成简单的文档排版。1984年，Leslie Lamport 在 \TeX{} 之上构建了 \LaTeX{}（读作 "lay-tech" 或 "lah-tech"）\cite{Lamport1994}，提供了一套面向文档逻辑结构的标记语言。用户只需声明文档类型、章节层级、图表标题等逻辑信息，\LaTeX{} 自动处理字体、间距、分页等视觉细节。

经过数十年的发展，\LaTeX{} 已形成庞大的宏包生态。截至2026年，CTAN（Comprehensive \TeX{} Archive Network）收录了超过6,000个宏包，覆盖了从化学分子式到音乐乐谱、从算法伪代码到语言学注释树的各种专业排版需求。

\subsection{\LaTeX{} 在学术写作中的意义}

相比于所见即所得（WYSIWYG）的文字处理软件，\LaTeX{} 在学术写作中具有以下核心优势：

\begin{enumerate}[label=\textbf{(\arabic*)}]
	\item \textbf{排版质量}：\TeX{} 的断行与分页算法基于动态规划全局优化，其数学公式排版至今仍是业界标杆。对于包含大量公式、交叉引用和浮动体的学位论文，\LaTeX{} 能自动维护编号一致性。
	\item \textbf{内容与样式分离}：作者专注于写作内容，排版格式由文档类与样式文件统一控制。同一份源文件可切换不同模板瞬间生成不同格式的输出。
	\item \textbf{纯文本与版本控制}：\texttt{.tex} 源码是纯文本文件，可以纳入 Git 等版本控制系统，实现协作写作、历史回溯与差异比较。
	\item \textbf{可复制性}：编译过程是确定性的，同一源码在不同机器上产生完全相同的 PDF。
\end{enumerate}

\section{\LaTeX{} 的基本操作}

\subsection{编译流程}

一个典型的 \LaTeX{} 文档编译流程如下：

\begin{enumerate}[label=\textbf{第\arabic*步：}]
	\item \textbf{XeLaTeX 第一遍}：解析源文件，生成 \texttt{.aux}（交叉引用信息）、\texttt{.toc}（目录信息）、\texttt{.out}（PDF书签）等辅助文件。
	\item \textbf{BibTeX}：读取 \texttt{.aux} 中的引用指令，从 \texttt{.bib} 数据库中提取被引条目，按指定格式生成 \texttt{.bbl}（参考文献排版数据）。
	\item \textbf{XeLaTeX 第二遍}：读入 \texttt{.bbl} 和上一遍的辅助文件，更新交叉引用。
	\item \textbf{XeLaTeX 第三遍}：解析上一遍更新的引用标签，确保所有编号稳定。
\end{enumerate}

可以使用 \texttt{latexmk} 工具自动化这一多遍编译流程。本模板根目录提供了 \texttt{.latexmkrc} 配置文件，执行 \texttt{latexmk -xelatex main} 即可一键编译。

\subsection{文档结构}

一个 \LaTeX{} 文档的基本结构为：

\begin{verbatim}
\documentclass{...}      % 文档类
% 导言区（加载宏包、定义命令）
\begin{document}
% 正文内容
\end{document}
\end{verbatim}

对于大型文档（如学位论文），通常将章节拆分到多个 \texttt{.tex} 文件中，由主文件通过 \verb|\include| 或 \verb|\input| 引入。本模板更进一步，使用 \texttt{subfiles} 宏包使每个子文件既能作为整体的一部分编译，也能独立编译预览。

\subsection{常用编译引擎}

\begin{table}[htbp]
	\centering
	\caption{\LaTeX{} 主流编译引擎对比}
	\label{tab:engines}
	\begin{tabular}{cccc}
		\toprule
		引擎 & 中文支持 & 字体支持 & 推荐度 \\
		\midrule
		pdfLaTeX & 需 CJK 宏包 & Type1/OT1 & \ding{55} \\
		XeLaTeX & 原生支持 & 系统字体 & \ding{51} \\
		LuaLaTeX & 原生支持 & 系统字体 & \ding{51} \\
		\bottomrule
	\end{tabular}
\end{table}

\textbf{本模板必须使用 XeLaTeX 编译}，原因是 \texttt{xeCJK} 宏包依赖于 XeTeX 引擎的原生 Unicode 和系统字体支持。

\section{为什么需要这个技术手册}

尽管 \LaTeX{} 功能强大，但其学习曲线相对陡峭：

\begin{itemize}
	\item \textbf{模块化带来的复杂性}：本模板将导言区拆分为 15 个独立的样式文件，虽然便于维护和 AI 协作，但普通用户初次面对时可能不知从何入手。
	\item \textbf{中文排版的特殊性}：中文字体配置、中文目录样式、GBT7714 参考文献格式等都是中文论文特有的需求，相关文档分散。
	\item \textbf{AI 协作是新范式}：本模板设计了 \texttt{CLAUDE.md} 和 \texttt{custom.tex} 来支持与 AI 助手的协作，这是一种新的工作方式，需要专门的说明。
	\item \textbf{学位论文的高要求}：湖南师范大学对学位论文的封面、扉页、摘要、目录、参考文献等有严格的格式规范，手动调整耗时且容易出错。
\end{itemize}

因此，本手册以模板自身作为排版引擎——你正在阅读的这份 PDF，正是由本模板编译生成的。手册中的每一个示例、每一段代码、每一个表格，都是"所见即所得"的模板效果展示。

\section{手册设计概述}

本手册按照学术论文章节结构组织，共分为七章和两个附录：

\begin{description}
	\item[第一章 绪论] 介绍 \TeX/\LaTeX{} 的历史背景、基本操作流程，以及本手册的编写目的和结构安排。
	\item[第二章 设计思路及结构] 详细解析模板的模块化设计理念、目录结构、各个样式文件的功能分工和用户变量系统。
	\item[第三章 使用说明] 提供从零开始的完整操作指南：填写个人信息、添加章节、插入图表公式、使用定理环境、管理参考文献以及切换学位类型。
	\item[第四章 AI 交互及 AI Agent 代理] 介绍如何利用 Claude Code 等 AI 助手与模板协作，\texttt{CLAUDE.md} 的机制，以及 \texttt{custom.tex} 的扩展工作流程。
	\item[第五章 更新及功能增加] 说明如何安全地给模板增加新功能——通过 \texttt{custom.tex} 添加宏包、定义命令、覆盖格式。
	\item[第六章 Git 提交] 介绍使用 Git 进行论文版本控制的最佳实践，包括 \texttt{.gitignore} 配置与协作策略。
	\item[第七章 结论] 总结最佳实践与推荐工作流程。
	\item[附录A] 提供 GB/T 7714 参考文献著录标准的常用条目格式示例。
	\item[附录B] 汇总常见 \LaTeX{} 编译错误及其解决方案，以及字体安装指南。
\end{description}

\end{document}
